\chapter{§3.10 锐角三角函数}
\label{sec:3.10}


\prerequisite{§3.4, §3.7}

\gradelevel{9 年级}


\section*{锐角三角函数的定义}


\begin{explore}


一座山的坡面与水平面的夹角是 $30°$ 。如果你沿着坡面走了 $100$ 米，你实际上升了多少米？

要回答这个问题，我们需要知道 $30°$ 角对应的"对边与斜边的比值"。这就是三角函数的核心思想——用比值来描述角与边的关系。

\end{explore}

\begin{understand}


在直角三角形中，设一个锐角为 $\angle A$ ，则：

\begin{figure}[htbp]
\centering
\begin{tikzpicture}[every node/.style={font=\small}, scale=1.0]
  \coordinate (C) at (0,0);
  \coordinate (A) at (4,0);
  \coordinate (B) at (4,3);
  \draw[main line] (A)--(B)--(C)--cycle;
  \node[below left]  at (C) {$C$};
  \node[below right] at (A) {$A$};
  \node[above right] at (B) {$B$};
  % Right angle at A
  \draw (3.75,0)--(3.75,0.25)--(4,0.25);
  % Angle α at C
  \draw[angle arc, draw=red!70!black, thick]
    (0:0.75) arc[start angle=0, end angle={atan2(3,4)}, radius=0.75];
  \node[red!70!black] at (18:1.1) {$\alpha$};
  % Side labels
  \node[below] at (2,0) {邻边（adjacent）};
  \node[right] at (4,1.5) {对边（opposite）};
  \node[above left, rotate={atan2(3,4)}] at (2,1.5) {斜边（hypotenuse）};
  % Trig ratios below
  \node[below=1.5cm] at (2,0)
    {$\sin\alpha = \dfrac{\text{对边}}{\text{斜边}},\quad
     \cos\alpha = \dfrac{\text{邻边}}{\text{斜边}},\quad
     \tan\alpha = \dfrac{\text{对边}}{\text{邻边}}$};
\end{tikzpicture}
\caption{锐角三角函数定义（以直角三角形中 $\angle C = \alpha$ 为例）}
\label{fig:trig-definition}
\end{figure}


\[
\sin A = \frac{\text{对边}}{\text{斜边}} = \frac{a}{c}
\]


\[
\cos A = \frac{\text{邻边}}{\text{斜边}} = \frac{b}{c}
\]


\[
\tan A = \frac{\text{对边}}{\text{邻边}} = \frac{a}{b}
\]


其中：
\begin{itemize}
  \item \textbf{对边}： $\angle A$ 的对面那条边（ $a$ ）
  \item \textbf{邻边}： $\angle A$ 的旁边那条直角边（ $b$ ）
  \item \textbf{斜边}：直角的对面那条边，即最长的边（ $c$ ）
\end{itemize}


\end{understand}

\begin{quote}
记忆口诀："正弦 $=$ 对比斜，余弦 $=$ 邻比斜，正切 $=$ 对比邻。"
\end{quote}


\textbf{为什么三角函数是确定的？}

在所有含 $\angle A$ 的直角三角形中，无论三角形的大小如何， $\sin A$ 、 $\cos A$ 、 $\tan A$ 的值都相同。这是因为这些直角三角形都是相似的（AA 相似：共同的 $\angle A$ 和 $90°$ 角），相似三角形的对应边成比例，因此比值不变。

\begin{workedexamples}


\textbf{例 1.} 如图，在 $Rt\triangle ABC$ 中， $\angle C = 90°$ ， $AC = 3$ ， $BC = 4$ ， $AB = 5$ 。求 $\sin A$ 、 $\cos A$ 、 $\tan A$ 。

\textbf{解：}
$\angle A$ 的对边 $= BC = 4$ ，邻边 $= AC = 3$ ，斜边 $= AB = 5$ 。

$\sin A = \dfrac{BC}{AB} = \dfrac{4}{5}$

$\cos A = \dfrac{AC}{AB} = \dfrac{3}{5}$

$\tan A = \dfrac{BC}{AC} = \dfrac{4}{3}$

\textbf{例 2.} 在 $Rt\triangle ABC$ 中， $\angle C = 90°$ ， $\sin A = \dfrac{5}{13}$ 。求 $\cos A$ 和 $\tan A$ 。

\textbf{解：}
设 $BC = 5k$ ， $AB = 13k$ （ $\sin A = \dfrac{BC}{AB} = \dfrac{5}{13}$ ）。

由勾股定理： $AC = \sqrt{AB^2 - BC^2} = \sqrt{169k^2 - 25k^2} = 12k$ 。

$\cos A = \dfrac{AC}{AB} = \dfrac{12k}{13k} = \dfrac{12}{13}$

$\tan A = \dfrac{BC}{AC} = \dfrac{5k}{12k} = \dfrac{5}{12}$

\textbf{例 3.} 在 $Rt\triangle ABC$ 中， $\angle C = 90°$ ， $\tan B = 2$ ， $BC = 6$ 。求 $AC$ 和 $AB$ 。

\textbf{解：}
$\tan B = \dfrac{AC}{BC} = \dfrac{AC}{6} = 2$ ， $AC = 12$ 。

$AB = \sqrt{BC^2 + AC^2} = \sqrt{36 + 144} = \sqrt{180} = 6\sqrt{5}$ 。

\end{workedexamples}

\begin{keytakeaway}


\begin{itemize}
  \item $\sin A = \dfrac{\text{对边}}{\text{斜边}}$ ， $\cos A = \dfrac{\text{邻边}}{\text{斜边}}$ ， $\tan A = \dfrac{\text{对边}}{\text{邻边}}$ 。
  \item 三角函数值只取决于角度大小，与三角形的大小无关。
  \item 已知一个三角函数值和一条边，可以求出其他边。
\end{itemize}


\end{keytakeaway}

\begin{practice}


\begin{enumerate}
  \item 在 $Rt\triangle ABC$ 中， $\angle C = 90°$ ， $AB = 10$ ， $AC = 8$ 。求 $\sin A$ 、 $\cos A$ 、 $\tan A$ 。
  \item 在 $Rt\triangle ABC$ 中， $\angle C = 90°$ ， $\cos B = \dfrac{3}{5}$ ， $AB = 15$ 。求 $BC$ 和 $AC$ 。
  \item 在 $Rt\triangle ABC$ 中， $\angle C = 90°$ ， $\tan A = \sqrt{3}$ 。求 $\angle A$ （利用下一节的特殊角值）。
\end{enumerate}


\end{practice}



\section*{特殊角的三角函数值}


\begin{explore}


虽然大多数角的三角函数值需要查表或用计算器，但有三个特殊角—— $30°$ 、 $45°$ 、 $60°$ ——的三角函数值可以精确计算出来，值得记住。

\end{explore}

\begin{understand}


\textbf{ $45°$ 角}：等腰直角三角形中，两直角边相等。设直角边为 $1$ ，则斜边为 $\sqrt{2}$ 。

\[
\sin 45° = \frac{1}{\sqrt{2}} = \frac{\sqrt{2}}{2}, \quad \cos 45° = \frac{\sqrt{2}}{2}, \quad \tan 45° = 1
\]


\textbf{ $30°$ 和 $60°$ 角}：将一个等边三角形沿高对折，得到含 $30°$ 和 $60°$ 角的直角三角形。设等边三角形边长为 $2$ ，则短直角边（ $30°$ 对边） $= 1$ ，长直角边（ $60°$ 对边） $= \sqrt{3}$ ，斜边 $= 2$ 。

\begin{figure}[htbp]
\centering
\begin{tikzpicture}[every node/.style={font=\small}, scale=0.9]
  % 30-60-90 triangle (left)
  \begin{scope}
    \coordinate (C1) at (0,0);
    \coordinate (A1) at (2,0);
    \coordinate (B1) at (2,{2*tan(60)});
    \draw[main line] (C1)--(A1)--(B1)--cycle;
    \draw (1.75,0)--(1.75,0.25)--(2,0.25);
    \draw[angle arc, draw=red!70!black]
      (0:0.5) arc[start angle=0,end angle=60,radius=0.5];
    \node[red!70!black] at (30:0.75) {$60°$};
    \draw[angle arc, draw=blue!70!black]
      ($(B1)+(270:0.5)$) arc[start angle=270,end angle={270-30},radius=0.5];
    \node[blue!70!black, font=\scriptsize] at ($(B1)+(-0.5,-0.5)$) {$30°$};
    \node[below] at (1,0) {$1$};
    \node[right] at (2,{tan(60)}) {$\sqrt{3}$};
    \node[above left, rotate=60] at (1,{tan(60)}) {$2$};
  \end{scope}
  % 45-45-90 triangle (right)
  \begin{scope}[xshift=5cm]
    \coordinate (C2) at (0,0);
    \coordinate (A2) at (2,0);
    \coordinate (B2) at (2,2);
    \draw[main line] (C2)--(A2)--(B2)--cycle;
    \draw (1.75,0)--(1.75,0.25)--(2,0.25);
    \draw[angle arc, draw=red!70!black]
      (0:0.5) arc[start angle=0,end angle=45,radius=0.5];
    \node[red!70!black] at (22:0.8) {$45°$};
    \node[below] at (1,0) {$1$};
    \node[right] at (2,1) {$1$};
    \node[above left, rotate=45] at (1,1) {$\sqrt{2}$};
  \end{scope}
\end{tikzpicture}
\caption{特殊角三角函数的精确值（$30°$-$60°$-$90°$ 和 $45°$-$45°$-$90°$ 三角形）}
\label{fig:special-angles-trig}
\end{figure}


\[
\sin 30° = \frac{1}{2}, \quad \cos 30° = \frac{\sqrt{3}}{2}, \quad \tan 30° = \frac{1}{\sqrt{3}} = \frac{\sqrt{3}}{3}
\]


\[
\sin 60° = \frac{\sqrt{3}}{2}, \quad \cos 60° = \frac{1}{2}, \quad \tan 60° = \sqrt{3}
\]


\textbf{特殊角三角函数值表}：

\begin{center}
\begin{tabular}{llll}
\toprule
角度 & $\sin$ & $\cos$ & $\tan$ \\
\midrule
$30°$ & $\dfrac{1}{2}$ & $\dfrac{\sqrt{3}}{2}$ & $\dfrac{\sqrt{3}}{3}$ \\
$45°$ & $\dfrac{\sqrt{2}}{2}$ & $\dfrac{\sqrt{2}}{2}$ & $1$ \\
$60°$ & $\dfrac{\sqrt{3}}{2}$ & $\dfrac{1}{2}$ & $\sqrt{3}$ \\
\bottomrule
\end{tabular}
\end{center}


\end{understand}

\begin{quote}
记忆技巧： $\sin$ 值从 $30°$ 到 $60°$ 依次增大： $\dfrac{1}{2} < \dfrac{\sqrt{2}}{2} < \dfrac{\sqrt{3}}{2}$ 。 $\cos$ 值的变化恰好相反。
\end{quote}


\begin{workedexamples}


\textbf{例 1.} 求 $\sin 30° + \cos 60°$ 的值。

\textbf{解：}
$\sin 30° + \cos 60° = \dfrac{1}{2} + \dfrac{1}{2} = 1$ 。

\textbf{例 2.} 求 $\sin^2 45° + \cos^2 45°$ 的值。

\textbf{解：}
$\sin^2 45° + \cos^2 45° = \left(\dfrac{\sqrt{2}}{2}\right)^2 + \left(\dfrac{\sqrt{2}}{2}\right)^2 = \dfrac{1}{2} + \dfrac{1}{2} = 1$ 。

\textbf{例 3.} 在 $Rt\triangle ABC$ 中， $\angle C = 90°$ ， $\angle A = 30°$ ， $AB = 10$ 。求 $BC$ 和 $AC$ 。

\textbf{解：}
$BC = AB \sin 30° = 10 \times \dfrac{1}{2} = 5$ 。

$AC = AB \cos 30° = 10 \times \dfrac{\sqrt{3}}{2} = 5\sqrt{3}$ 。

验证： $BC^2 + AC^2 = 25 + 75 = 100 = AB^2$ 。正确。

\end{workedexamples}

\begin{keytakeaway}


\begin{itemize}
  \item $30°$ 、 $45°$ 、 $60°$ 的三角函数值是精确值，必须记住。
  \item $\sin$ 从 $30°$ 到 $60°$ 递增， $\cos$ 递减。
  \item $\tan 45° = 1$ 是一个重要的分界点。
\end{itemize}


\end{keytakeaway}

\begin{practice}


\begin{enumerate}
  \item 计算 $2\sin 30° \times \cos 60° + \tan 45°$ 。
  \item 在 $Rt\triangle ABC$ 中， $\angle C = 90°$ ， $\angle B = 60°$ ， $AC = 6$ 。求 $AB$ 和 $BC$ 。
  \item 已知 $\sin \alpha = \dfrac{\sqrt{3}}{2}$ ， $\alpha$ 为锐角，求 $\alpha$ 。
\end{enumerate}


\end{practice}



\section*{三角函数之间的关系}


\begin{explore}


观察特殊角的三角函数值表，你能发现一些规律吗？比如 $\sin 30° = \cos 60°$ ——一个角的正弦等于另一个角（它的余角）的余弦。这不是巧合！

\end{explore}

\begin{understand}


设直角三角形中两个锐角分别为 $\angle A$ 和 $\angle B$ （ $\angle A + \angle B = 90°$ ）。

\textbf{互余关系}：

\[
\sin A = \cos B = \cos(90° - A)
\]


\[
\cos A = \sin B = \sin(90° - A)
\]


即：一个锐角的正弦等于其余角的余弦，反之亦然。

\end{understand}

\begin{quote}
这就是"余弦"名称的由来——"余角的正弦"。
\end{quote}


\textbf{商的关系}：

\[
\tan A = \frac{\sin A}{\cos A}
\]


因为 $\tan A = \dfrac{a}{b}$ ，而 $\sin A = \dfrac{a}{c}$ ， $\cos A = \dfrac{b}{c}$ ，所以 $\dfrac{\sin A}{\cos A} = \dfrac{a/c}{b/c} = \dfrac{a}{b} = \tan A$ 。

\textbf{平方关系}：

\[
\sin^2 A + \cos^2 A = 1
\]


因为 $\sin^2 A + \cos^2 A = \dfrac{a^2}{c^2} + \dfrac{b^2}{c^2} = \dfrac{a^2 + b^2}{c^2} = \dfrac{c^2}{c^2} = 1$ （勾股定理）。

\begin{workedexamples}


\textbf{例 1.} 已知 $\sin 28° \approx 0.469$ ，求 $\cos 62°$ 。

\textbf{解：}
因为 $28° + 62° = 90°$ ，所以 $\cos 62° = \sin 28° \approx 0.469$ 。

\textbf{例 2.} 已知 $\sin \alpha = \dfrac{3}{5}$ （ $\alpha$ 为锐角），求 $\cos \alpha$ 和 $\tan \alpha$ 。

\textbf{解：}
由 $\sin^2 \alpha + \cos^2 \alpha = 1$ ：

$\cos^2 \alpha = 1 - \sin^2 \alpha = 1 - \dfrac{9}{25} = \dfrac{16}{25}$ 。

因为 $\alpha$ 是锐角， $\cos \alpha > 0$ ，所以 $\cos \alpha = \dfrac{4}{5}$ 。

$\tan \alpha = \dfrac{\sin \alpha}{\cos \alpha} = \dfrac{3/5}{4/5} = \dfrac{3}{4}$ 。

\textbf{例 3.} 化简 $\dfrac{\sin 50°}{\cos 40°}$ 。

\textbf{解：}
$\cos 40° = \sin(90° - 40°) = \sin 50°$ 。

$\dfrac{\sin 50°}{\cos 40°} = \dfrac{\sin 50°}{\sin 50°} = 1$ 。

\textbf{例 4.} 化简 $\sin^2 25° + \sin^2 65°$ 。

\textbf{解：}
$\sin 65° = \cos 25°$ （互余）。

$\sin^2 25° + \sin^2 65° = \sin^2 25° + \cos^2 25° = 1$ 。

\end{workedexamples}

\begin{keytakeaway}


\begin{itemize}
  \item $\sin A = \cos(90° - A)$ ， $\cos A = \sin(90° - A)$ 。
  \item $\tan A = \dfrac{\sin A}{\cos A}$ 。
  \item $\sin^2 A + \cos^2 A = 1$ 。
\end{itemize}


\end{keytakeaway}

\begin{practice}


\begin{enumerate}
  \item 已知 $\cos 35° \approx 0.819$ ，求 $\sin 55°$ 。
  \item 已知 $\cos \alpha = \dfrac{5}{13}$ ，求 $\sin \alpha$ 和 $\tan \alpha$ 。
  \item 化简 $\cos^2 40° + \cos^2 50°$ 。
\end{enumerate}


\end{practice}



\section*{解直角三角形}


\begin{explore}


消防员要用云梯救人。已知窗户离地面 $12$ 米，云梯与地面的夹角不能超过 $60°$ （否则不安全）。云梯至少需要多长？

这就是"解直角三角形"——已知直角三角形的部分边和角，求出其余的边和角。

\end{explore}

\begin{understand}


\textbf{解直角三角形}就是在直角三角形中，根据已知的边和角，求出所有未知的边和角。

在 $Rt\triangle ABC$ （ $\angle C = 90°$ ）中，已知条件至少包含\textbf{一条边}和\textbf{另一个元素}（一条边或一个锐角）。

常用关系：
\begin{itemize}
  \item $\angle A + \angle B = 90°$
  \item $a^2 + b^2 = c^2$ （勾股定理）
  \item $\sin A = \dfrac{a}{c}$ ， $\cos A = \dfrac{b}{c}$ ， $\tan A = \dfrac{a}{b}$
\end{itemize}


\end{understand}

\begin{workedexamples}


\textbf{例 1.} 在 $Rt\triangle ABC$ 中， $\angle C = 90°$ ， $\angle A = 35°$ ， $c = 20$ 。求 $a$ 和 $b$ 。（ $\sin 35° \approx 0.574$ ， $\cos 35° \approx 0.819$ ）

\textbf{解：}
$\angle B = 90° - 35° = 55°$ 。

$a = c \sin A = 20 \times 0.574 = 11.48$ 。

$b = c \cos A = 20 \times 0.819 = 16.38$ 。

验证： $a^2 + b^2 \approx 131.8 + 268.3 = 400.1 \approx 400 = c^2$ 。基本正确（误差来自近似值）。

\textbf{例 2.} 在 $Rt\triangle ABC$ 中， $\angle C = 90°$ ， $a = 5$ ， $b = 8$ 。求 $\angle A$ 和 $c$ 。

\textbf{解：}
$\tan A = \dfrac{a}{b} = \dfrac{5}{8} = 0.625$ 。

查表或用计算器得 $\angle A \approx 32°$ 。

$\angle B = 90° - 32° = 58°$ 。

$c = \sqrt{a^2 + b^2} = \sqrt{25 + 64} = \sqrt{89} \approx 9.43$ 。

\textbf{例 3.} 消防云梯问题：窗户离地面 $12$ m，云梯与地面的夹角为 $60°$ 。求云梯的长度和梯脚到墙壁的距离。

\textbf{解：}
设云梯长为 $l$ ，梯脚到墙壁的距离为 $d$ 。

$\sin 60° = \dfrac{12}{l}$ ， $\dfrac{\sqrt{3}}{2} = \dfrac{12}{l}$ ， $l = \dfrac{24}{\sqrt{3}} = \dfrac{24\sqrt{3}}{3} = 8\sqrt{3} \approx 13.86$ m。

$\cos 60° = \dfrac{d}{l}$ ， $d = l \cos 60° = 8\sqrt{3} \times \dfrac{1}{2} = 4\sqrt{3} \approx 6.93$ m。

或直接： $\tan 60° = \dfrac{12}{d}$ ， $d = \dfrac{12}{\sqrt{3}} = 4\sqrt{3}$ m。

\end{workedexamples}

\begin{keytakeaway}


\begin{itemize}
  \item 解直角三角形需要至少两个已知量（其中至少一条边）。
  \item 根据已知和未知的组合，选择合适的三角函数关系。
  \item 解题时优先使用包含"一个已知量和一个未知量"的等式。
\end{itemize}


\end{keytakeaway}

\begin{practice}


\begin{enumerate}
  \item 在 $Rt\triangle ABC$ 中， $\angle C = 90°$ ， $\angle A = 45°$ ， $a = 7$ 。求 $b$ 和 $c$ 。
  \item 在 $Rt\triangle ABC$ 中， $\angle C = 90°$ ， $a = 8$ ， $c = 10$ 。求 $\angle A$ 、 $\angle B$ 和 $b$ 。
  \item 一根电线杆高 $10$ m，从杆顶向地面拉一根绳子，绳子与地面成 $30°$ 角。求绳子的长度和绳子固定点到杆底的距离。
\end{enumerate}


\end{practice}



\section*{应用：测量、坡度、仰角与俯角}


\begin{explore}


古人没有飞机和卫星，却能测量金字塔的高度、河流的宽度。他们的秘密武器就是三角函数。今天，测量员、建筑师、飞行员依然在日常工作中使用三角函数。

\end{explore}

\begin{understand}


\textbf{仰角与俯角}：

从水平线向上看某个目标的角度叫做\textbf{仰角}（elevation angle）。

从水平线向下看某个目标的角度叫做\textbf{俯角}（depression angle）。

\begin{figure}[htbp]
\centering
\begin{tikzpicture}[every node/.style={font=\small}]
  % Observer at O
  \coordinate (O) at (0,0);
  \fill (O) circle (2pt) node[left] {$O$（观测者）};
  % Horizontal line
  \draw[ext line] (-0.5,0) -- (4.5,0);
  % Elevation angle: target T above
  \coordinate (T) at (4,2.5);
  \fill (T) circle (2pt) node[right] {目标 $T$（在上方）};
  \draw[main line] (O)--(T);
  \draw[angle arc, draw=red!70!black, thick]
    (0:1.2) arc[start angle=0,end angle={atan2(2.5,4)},radius=1.2];
  \node[red!70!black] at ({atan2(2.5,4)/2}:1.6) {仰角 $\theta$};
  % Depression angle: target D below
  \coordinate (D) at (4,-1.8);
  \fill (D) circle (2pt) node[right] {目标 $D$（在下方）};
  \draw[main line] (O)--(D);
  \draw[angle arc, draw=blue!70!black, thick]
    (0:1.2) arc[start angle=0,end angle={-atan2(1.8,4)},radius=1.2];
  \node[blue!70!black] at ({-atan2(1.8,4)/2}:1.7) {俯角 $\phi$};
\end{tikzpicture}
\caption{仰角与俯角：从水平视线分别向上（仰角）和向下（俯角）量起}
\label{fig:elevation-depression}
\end{figure}


\textbf{坡度}：

坡面的铅直高度与水平距离的比叫做\textbf{坡度}（也叫坡比）。

\[
\text{坡度} = \tan \alpha = \frac{\text{铅直高度}}{\text{水平距离}}
\]


其中 $\alpha$ 是坡面与水平面的夹角（坡角）。

坡度常用比值或百分数表示：坡度 $1:2$ 表示水平前进 $2$ m 上升 $1$ m；坡度 $6\%$ 表示水平前进 $100$ m 上升 $6$ m。

\end{understand}

\begin{workedexamples}


\textbf{例 1.} 小明站在距一栋楼 $50$ m 远的地面上，测得楼顶的仰角为 $45°$ 。小明的眼睛离地面 $1.5$ m。求楼的高度。

\textbf{解：}
设楼顶到小明眼睛水平线以上的高度为 $h$ 。

$\tan 45° = \dfrac{h}{50}$ ， $1 = \dfrac{h}{50}$ ， $h = 50$ m。

楼的高度 $= h + 1.5 = 50 + 1.5 = 51.5$ m。

\textbf{例 2.} 从山顶测得山脚处一棵树底部的俯角为 $60°$ ，树顶的俯角为 $45°$ 。山高 $100$ m。求树的高度。

\begin{figure}[htbp]
\centering
\begin{tikzpicture}[every node/.style={font=\small}, scale=0.85]
  % Building
  \draw[main line] (0,0) -- (0,4) -- (0.5,4) -- (0.5,0) -- cycle;
  \fill[gray!20] (0,0) -- (0,4) -- (0.5,4) -- (0.5,0) -- cycle;
  % Ground
  \draw[main line] (-0.5,0) -- (5,0);
  % Observer O at top of building
  \coordinate (O) at (0.25,4);
  \fill (O) circle (2pt) node[above left] {$O$};
  % Target A on ground
  \coordinate (A) at (4,0);
  \fill (A) circle (2pt) node[below] {$A$};
  % Horizontal line from O
  \draw[ext line] (-0.5,4) -- (5,4);
  % Line OA
  \draw[main line] (O)--(A);
  % Depression angle
  \draw[angle arc, draw=red!70!black, thick]
    ($(O)+(0:0.8)$) arc[start angle=0,end angle={-atan2(4,3.75)},radius=0.8];
  \node[red!70!black, font=\scriptsize] at ($(O)+({-atan2(4,3.75)/2}:1.2)$) {$\alpha$（俯角）};
  % Height label
  \draw[<->, >=Stealth, thin] (-0.7,0) -- (-0.7,4) node[midway,left] {$h$};
  % Horizontal distance
  \draw[<->, >=Stealth, thin] (0.25,-0.5) -- (4,-0.5) node[midway,below] {$d$};
  \node[font=\scriptsize] at (2.5,2) {$\tan\alpha = \dfrac{h}{d}$};
\end{tikzpicture}
\caption{俯角例题：$\tan\alpha = h/d$，已知两量可求第三量}
\label{fig:depression-angle-ex}
\end{figure}


\textbf{解：}
设山脚到树的水平距离为 $d$ ，树高为 $h$ 。

从山顶到树底部的俯角 $60°$ ： $\tan 60° = \dfrac{100}{d}$ ， $d = \dfrac{100}{\sqrt{3}} = \dfrac{100\sqrt{3}}{3}$ m。

从山顶到树顶的俯角 $45°$ ： $\tan 45° = \dfrac{100 - h}{d}$ （树顶比山顶低 $100 - h$ ）。

$1 = \dfrac{100 - h}{\dfrac{100\sqrt{3}}{3}}$

$100 - h = \dfrac{100\sqrt{3}}{3}$

$h = 100 - \dfrac{100\sqrt{3}}{3} = \dfrac{300 - 100\sqrt{3}}{3} = \dfrac{100(3 - \sqrt{3})}{3} \approx \dfrac{100 \times 1.268}{3} \approx 42.3$ m。

\textbf{例 3.} 一段公路的坡度为 $1:5$ ，坡长 $200$ m。求这段公路的铅直高度和水平距离。

\textbf{解：}
坡度 $= \tan \alpha = \dfrac{1}{5}$ 。

设铅直高度为 $h$ ，水平距离为 $d$ 。 $\dfrac{h}{d} = \dfrac{1}{5}$ ， $h = \dfrac{d}{5}$ 。

坡长（斜面长） $= \sqrt{h^2 + d^2} = 200$ 。

$\sqrt{\dfrac{d^2}{25} + d^2} = 200$

$\sqrt{\dfrac{26d^2}{25}} = 200$

$\dfrac{d\sqrt{26}}{5} = 200$

$d = \dfrac{1000}{\sqrt{26}} \approx 196.1$ m。

$h = \dfrac{d}{5} \approx 39.2$ m。

\textbf{例 4.} 要测量河对岸一棵树 $A$ 到河岸的距离。在河这边的岸上取两点 $B$ 、 $C$ ， $BC = 100$ m， $\angle ABC = 90°$ ， $\angle ACB = 60°$ 。求 $AB$ 的长。

\textbf{解：}
在 $Rt\triangle ABC$ （ $\angle B = 90°$ ）中：

$\tan C = \dfrac{AB}{BC}$

$\tan 60° = \dfrac{AB}{100}$

$\sqrt{3} = \dfrac{AB}{100}$

$AB = 100\sqrt{3} \approx 173.2$ m。

\end{workedexamples}

\begin{keytakeaway}


\begin{itemize}
  \item 仰角：向上看的角度；俯角：向下看的角度。
  \item 坡度 $= \tan(\text{坡角}) = \dfrac{\text{铅直高度}}{\text{水平距离}}$ 。
  \item 实际测量问题的关键：识别直角三角形，确定已知和未知元素。
\end{itemize}


\end{keytakeaway}

\begin{practice}


\begin{enumerate}
  \item 小华站在距塔底 $30$ m 处，测得塔顶仰角为 $60°$ 。小华身高忽略不计，求塔高。
  \item 一段坡路的坡角为 $30°$ ，沿坡面走了 $600$ m，实际升高了多少米？水平前进了多少米？
  \item 飞机在 $1000$ m 高空飞行，飞行员看到前方地面目标的俯角为 $45°$ 。求飞机到目标的水平距离。
  \item 在河岸一侧取 $B$ 、 $C$ 两点， $BC = 80$ m， $\angle B = 90°$ ，测得对岸标志物 $A$ 使 $\angle ACB = 45°$ 。求 $AB$ 。
\end{enumerate}


\end{practice}



\begin{answer}


\begin{enumerate}
  \item $BC = \sqrt{AB^2 - AC^2} = \sqrt{100 - 64} = 6$ 。 $\sin A = \dfrac{6}{10} = \dfrac{3}{5}$ ， $\cos A = \dfrac{8}{10} = \dfrac{4}{5}$ ， $\tan A = \dfrac{6}{8} = \dfrac{3}{4}$ 。
\end{enumerate}


\begin{enumerate}
  \item $\cos B = \dfrac{BC}{AB} = \dfrac{3}{5}$ ， $BC = AB \cos B = 15 \times \dfrac{3}{5} = 9$ 。 $AC = \sqrt{AB^2 - BC^2} = \sqrt{225 - 81} = \sqrt{144} = 12$ 。
\end{enumerate}


\begin{enumerate}
  \item $\tan A = \sqrt{3} = \tan 60°$ ，所以 $\angle A = 60°$ 。
\end{enumerate}


\begin{enumerate}
  \item $2 \times \dfrac{1}{2} \times \dfrac{1}{2} + 1 = \dfrac{1}{2} + 1 = \dfrac{3}{2}$ 。
\end{enumerate}


\begin{enumerate}
  \item $\angle A = 90° - 60° = 30°$ 。 $\tan B = \dfrac{AC}{BC}$ ， $\tan 60° = \dfrac{6}{BC}$ ， $BC = \dfrac{6}{\sqrt{3}} = 2\sqrt{3}$ 。 $AB = \dfrac{AC}{\sin B} = \dfrac{6}{\sin 60°} = \dfrac{6}{\sqrt{3}/2} = 4\sqrt{3}$ 。
\end{enumerate}


\begin{enumerate}
  \item $\sin \alpha = \dfrac{\sqrt{3}}{2} = \sin 60°$ ， $\alpha = 60°$ 。
\end{enumerate}


\begin{enumerate}
  \item $\sin 55° = \cos 35° \approx 0.819$ 。
\end{enumerate}


\begin{enumerate}
  \item $\sin \alpha = \sqrt{1 - \cos^2 \alpha} = \sqrt{1 - \dfrac{25}{169}} = \sqrt{\dfrac{144}{169}} = \dfrac{12}{13}$ 。 $\tan \alpha = \dfrac{12/13}{5/13} = \dfrac{12}{5}$ 。
\end{enumerate}


\begin{enumerate}
  \item $\cos^2 40° + \cos^2 50° = \cos^2 40° + \sin^2 40° = 1$ （因为 $\cos 50° = \sin 40°$ ）。
\end{enumerate}


\begin{enumerate}
  \item $\angle B = 45°$ 。 $\tan A = \dfrac{a}{b}$ ， $\tan 45° = \dfrac{7}{b}$ ， $b = 7$ 。 $c = \sqrt{49 + 49} = 7\sqrt{2}$ 。
\end{enumerate}


\begin{enumerate}
  \item $b = \sqrt{c^2 - a^2} = \sqrt{100 - 64} = 6$ 。 $\sin A = \dfrac{8}{10} = \dfrac{4}{5} = 0.8$ ， $\angle A \approx 53.13°$ （或查表 $\angle A \approx 53°$ ）。 $\angle B \approx 37°$ 。
\end{enumerate}


\begin{enumerate}
  \item $\sin 30° = \dfrac{10}{l}$ ， $l = \dfrac{10}{\sin 30°} = 20$ m。 $d = 10 \times \dfrac{1}{\tan 30°} = 10\sqrt{3} \approx 17.32$ m。
\end{enumerate}


\begin{enumerate}
  \item $\tan 60° = \dfrac{h}{30}$ ， $h = 30\sqrt{3} \approx 51.96$ m。
\end{enumerate}


\begin{enumerate}
  \item 升高 $= 600 \times \sin 30° = 600 \times \dfrac{1}{2} = 300$ m。水平前进 $= 600 \times \cos 30° = 600 \times \dfrac{\sqrt{3}}{2} = 300\sqrt{3} \approx 519.6$ m。
\end{enumerate}


\begin{enumerate}
  \item $\tan 45° = \dfrac{1000}{d}$ ， $d = 1000$ m。
\end{enumerate}


\begin{enumerate}
  \item $\tan 45° = \dfrac{AB}{BC}$ ， $AB = BC \tan 45° = 80 \times 1 = 80$ m。
\end{enumerate}


\end{answer}
